\documentclass[11pt,a4paper]{article}

\usepackage[utf8]{inputenc}
\usepackage[T1]{fontenc}
\usepackage{lmodern}
\usepackage[margin=1in]{geometry}
\usepackage{amsmath, amssymb, amsthm}
\usepackage{microtype}
\usepackage{enumitem}
\usepackage{xcolor}
\usepackage{hyperref}

\hypersetup{
    colorlinks=true,
    linkcolor=blue!50!black,
    urlcolor=blue!50!black,
    citecolor=blue!50!black,
}

\newcommand{\R}{\mathbb{R}}
\newcommand{\N}{\mathbb{N}}
\newcommand{\E}{\mathbb{E}}
\newcommand{\X}{\mathcal{X}}
\newcommand{\A}{\mathcal{A}}
\newcommand{\F}{\mathcal{F}}
\newcommand{\Tdisc}{T_{\mathrm{disc}}}
\newcommand{\Tenv}{T_{\mathrm{env}}}
\newcommand{\Tca}{T_{\mathrm{CA}(1)}}
\newcommand{\Td}{T_{\mathrm{D1}(1)}}
\newcommand{\Cx}{CX}
\newcommand{\Iclass}{I}

\newtheorem{lemma}{Lemma}
\newtheorem{theorem}{Theorem}
\newtheorem{proposition}{Proposition}
\theoremstyle{definition}
\newtheorem{remark}{Remark}

\title{Threshold Optimality for an MMPP/M/1 Queue with Admission Control\\[0.3em]
       \large Homework Set 2, Problem 1.3}
\author{Leonardo Cruciani}
\date{May 2026}

\begin{document}
\maketitle

\section{The model}

We consider a single-server queue with Markov-modulated Poisson arrivals,
exponential service, and admission control. The state is
\[
    (x, y) \in \N_0 \times \{1, 2\},
\]
where $x$ is the number of customers in the system and $y$ is the
state of the modulator. In modulator state $y$, customers arrive at
rate $\lambda_y$ with $\lambda_1 < \lambda_2$. The modulator switches
to the other state at rate $q$, independently of the queue. The
server processes at rate $\mu$. An arriving customer can be rejected
at cost $R$ per rejection; a holding cost $c$ is paid per unit time
per customer in the system. With discount factor $\alpha \in (0, 1)$
and after uniformisation at rate $\gamma = \lambda_2 + \mu + q$, the
$n$-stage value function satisfies, for all $(x, y) \in \N_0 \times
\{1, 2\}$,
\begin{equation}
\label{eq:V}
\begin{aligned}
V_n(x, y) \;=\; cx \;+\; \alpha\Big[\;
        & \lambda_y\, \min\{R + V_{n-1}(x, y),\, V_{n-1}(x + 1, y)\} \\
    +\, & \mu\, V_{n-1}\big((x - 1)^+,\, y\big) \\
    +\, & q\, V_{n-1}(x, y') \\
    +\, & (\lambda_2 - \lambda_y)\, V_{n-1}(x, y)
\;\Big],
\end{aligned}
\end{equation}
with $V_0 \equiv 0$ and $y' = 3 - y$. The terms inside the bracket
correspond, in order, to a (possibly rejected) admission, a potential
service completion, a modulator transition, and a dummy self-loop
introduced by uniformisation.

The aim of this report is to prove

\begin{theorem}[Threshold optimality]
\label{thm:5}
There exist thresholds $U_1, U_2 \in \N_0$ such that, for the
discounted cost criterion, it is optimal to reject an arriving
customer in state $(x, y)$ if and only if $x \ge U_y$.
\end{theorem}

The proof proceeds by induction on $n$ in the event-based dynamic
programming framework of Koole (2006); references of the form
``Thm.~7.\textit{n}'' or ``Def.~5.\textit{n}'' below refer to that
framework. The strategy is:

\begin{enumerate}[leftmargin=2em, itemsep=0pt]
    \item Decompose~\eqref{eq:V} into event operators (\S\ref{sec:decomp}).
    \item Identify the functional class to propagate
        (\S\ref{sec:class}).
    \item Verify by induction that each operator preserves that class
        (\S\ref{sec:prop}).
    \item Conclude threshold optimality by control monotonicity
        (\S\ref{sec:concl}).
\end{enumerate}

A brief remark on the optional sharper statement $U_1 \ge U_2$ closes
the report (\S\ref{sec:thm6}).

\section{Operator decomposition}
\label{sec:decomp}

We identify the value-function components with operators from the
framework's catalogue (Defs.~5.1--5.3). The state-space convention
treats $y$ as the environment dimension (dimension $0$) and $x$ as
the single queue dimension (dimension $1$).

\paragraph{Discount operator.}
The direct cost and discount, by Def.~4.1, is
\[
    \Tdisc f(x, y) \;=\; cx \,+\, \alpha\, f(x, y).
\]

\paragraph{Controlled-arrival operator.}
The admission decision is, by Def.~5.2,
\[
    \Tca f(x, y) \;=\; \min\{R + f(x, y),\, f(x + 1, y)\},
\]
i.e.\ $\Tca$ with rejection cost $R$ and arrival cost $0$.

\paragraph{Single-server departure operator.}
The service event is, by Def.~5.3,
\[
    \Td f(x, y) \;=\; f\big((x - 1)^+, y\big).
\]

\paragraph{Environment operator.}
The mixing of the four events (admission, departure, modulator-jump,
dummy self-loop) and the $y \to y'$ transition of the modulator are
both bundled into the environment operator of Def.~4.2, which in
Koole's notation reads
\begin{equation}
\label{eq:Tenv-Koole}
\Tenv(F_1, \ldots, F_l)(X) \;=\; \sum_{Y \in \N_0} \Lambda(X_0, Y)\,
\sum_{j=1}^{l} Q^j(X_0, Y)\, F_j(X^*),
\end{equation}
where $X$ denotes the full state, $X_0$ is the environment coordinate,
$Y$ is the dummy summation variable ranging over possible new
environment values, $X^*$ is $X$ with the 0th coordinate replaced by
$Y$, $\Lambda(X_0, Y)$ is the probability the environment moves from
$X_0$ to $Y$, and $Q^j(X_0, Y)$ is the conditional probability that
mark $j$ (i.e.\ which of the slots $F_1, \ldots, F_l$ fires) occurs
given that transition. (Capitals denote Koole's abstract symbols;
when we instantiate below for the MMPP, the env coordinate $X_0$
becomes our $y \in \{1, 2\}$, the dummy summation $Y$ ranges over
$\{y, y'\}$, and so on.) In the MMPP instance we have $l = 4$
slots:
\begin{itemize}[leftmargin=2em, itemsep=2pt]
    \item $f_1$ admission, $f_2$ departure, $f_3$ modulator-jump,
        $f_4$ dummy self-loop;
    \item the environment kernel is
        $\lambda(y, y) = (\lambda_2 + \mu)/\gamma$ (env stays) and
        $\lambda(y, y') = q/\gamma$ (env jumps), with
        $\gamma = \lambda_2 + \mu + q$ the uniformisation rate;
    \item when the env \emph{stays}, the firing mark is admission,
        departure, or dummy, with conditional probabilities
        $q^1(y, y) = \lambda_y/(\lambda_2 + \mu)$,
        $q^2(y, y) = \mu/(\lambda_2 + \mu)$,
        $q^4(y, y) = (\lambda_2 - \lambda_y)/(\lambda_2 + \mu)$,
        and $q^3(y, y) = 0$;
    \item when the env \emph{jumps}, only the modulator-jump mark
        fires: $q^3(y, y') = 1$, others zero.
\end{itemize}
Substituting these into \eqref{eq:Tenv-Koole} and collecting terms
gives the expanded form
\begin{equation}
\label{eq:Tenv-expanded}
\Tenv(f_1, f_2, f_3, f_4)(x, y) =
\tfrac{1}{\gamma}\!\left[
    \lambda_y\, f_1(x, y)
    + \mu\, f_2(x, y)
    + q\, \underbrace{f_3(x, y')}_{\text{via Koole's }X^*}
    + (\lambda_2 - \lambda_y)\, f_4(x, y)
\right].
\end{equation}
The crucial point is that the $y \to y'$ relabeling of $f_3$ does
\emph{not} require a separate operator: it is produced by the $X^*$
mechanism inside Def.~4.2, which evaluates each slot at the new
environment coordinate selected by the outer sum.

\paragraph{Putting it together.}
Equation~\eqref{eq:V} can be rewritten in pure Koole form as
\begin{equation}
\label{eq:Vop}
V_n \;=\; \Tdisc \Bigl(\, \alpha' \, \Tenv\bigl[\Tca V_{n-1},\, \Td V_{n-1},\, V_{n-1},\, V_{n-1}\bigr] \Bigr),
\end{equation}
where $\alpha' = \alpha \gamma$ absorbs the rate factor (the bracket
in~\eqref{eq:V} is exactly $\gamma$ times the convex combination
in~\eqref{eq:Tenv-expanded}). Note that the third $V_{n-1}$ argument
appears in operator notation \emph{without any $y'$}: the
$V_{n-1}(x, y')$ term of \eqref{eq:V} is recovered by the $X^*$
substitution inside $\Tenv$ when the environment jumps. Concretely,
expanding \eqref{eq:Vop} via \eqref{eq:Tenv-expanded} reproduces
\eqref{eq:V} term by term.

\section{The functional class}
\label{sec:class}

By Theorem~3.1 of the framework, if a class $\F$ of functions
contains $V_0$ and is preserved by every operator appearing in
\eqref{eq:Vop}, then $V_n \in \F$ for all $n$. We take
\[
    \F \;=\; \Iclass(1) \cap \Cx(1),
\]
the class of functions $f : \N_0^2 \to \R$ that are
\emph{non-decreasing in $x$ for each $y$}
(Def.~6.1, $f \in I(1)$) and \emph{convex in $x$ for each $y$}
(Def.~6.3, $f \in \Cx(1)$):
\begin{align*}
    f \in \Iclass(1) :\;&\;
        f(x, y) \le f(x + 1, y) \quad \forall (x, y) \in \N_0^2,\\
    f \in \Cx(1)     :\;&\;
        2 f(x + 1, y) \le f(x, y) + f(x + 2, y) \quad \forall (x, y) \in \N_0^2.
\end{align*}
The zero function $V_0 \equiv 0$ trivially lies in $\F$.

The choice of $\F$ is motivated by Theorem~8.1 of the framework
(stated for $\Tca$ here): if $f \in \Cx(1)$ then the minimising action
in $\Tca f$ is non-increasing in $x$. Combined with finiteness of the
action set, this gives a threshold structure in $x$ for each $y$,
which is exactly the statement of Theorem~\ref{thm:5}.

\section{Propagation lemmas}
\label{sec:prop}

We show $\F = \Iclass(1) \cap \Cx(1)$ is preserved by every operator
in~\eqref{eq:Vop}.

\begin{lemma}[Increasingness propagation]
\label{lem:I}
If $f_1, \ldots, f_4 \in \Iclass(1)$, then $\Tdisc f_1$, $\Tca f_1$,
$\Td f_1$, and $\Tenv(f_1, f_2, f_3, f_4)$ are all in $\Iclass(1)$.
\end{lemma}

\begin{proof}
\textbf{(a)} $\Tdisc$. We compute
\[
    \Tdisc f(x, y) - \Tdisc f(x - 1, y)
    \;=\; c + \alpha\, \bigl[f(x, y) - f(x - 1, y)\bigr]
    \;\ge\; 0,
\]
since $c > 0$, $\alpha > 0$, and $f \in \Iclass(1)$. Hence $\Tdisc f
\in \Iclass(1)$. (This is the $\Tdisc$ row of Thm.~7.1 with
$C(x, y) = cx \in \Iclass(1)$.)

\textbf{(b)} $\Tca$. By Thm.~7.2 of the framework, $\Tca : \Iclass(1)
\to \Iclass(1)$ (the controlled-arrival operator preserves
increasingness in the controlled dimension; here, in $x$).

\textbf{(c)} $\Td$. By Thm.~7.3, $\Td : \Iclass(1) \to \Iclass(1)$
(the single-server departure operator preserves increasingness).

\textbf{(d)} $\Tenv$. By Def.~4.2,
$\Tenv(f_1, \ldots, f_4)(x, y)$ is a convex combination of the
values $f_j(x, y_{\mathrm{new}})$ for $y_{\mathrm{new}} \in \{y, y'\}$
and $j \in \{1, 2, 3, 4\}$, with non-negative weights
$\Lambda(y, y_{\mathrm{new}})\, Q^j(y, y_{\mathrm{new}})$ (Koole's
environment kernel and mark probabilities) that sum to $1$. Each
$f_j$ is in $\Iclass(1)$ by assumption, meaning each $y$-slice
$f_j(\cdot, y_{\mathrm{new}})$ is non-decreasing in $x$. Closure of
$\Iclass(1)$ under non-negative convex combinations gives the
result. This is the $\Tenv$ row of Thm.~7.1, which is what makes the
environment dimension ``free'' for propagation: as long as the
environment kernel $\Lambda$ does not depend on the queue coordinate
$x$ (here it does not --- our $q$ is independent of $x$), $\Tenv$
preserves $\Iclass(1)$.
\end{proof}

\begin{lemma}[Convexity propagation]
\label{lem:Cx}
If $f_1, \ldots, f_4 \in \Iclass(1) \cap \Cx(1)$, then $\Tdisc f_1$,
$\Tca f_1$, $\Td f_1$, and $\Tenv(f_1, f_2, f_3, f_4)$ are all in
$\Iclass(1) \cap \Cx(1)$.
\end{lemma}

\begin{proof}
\textbf{(a)} $\Tdisc$. Convexity in $x$:
\begin{align*}
    \Tdisc f(x, y) + \Tdisc f(x + 2, y) - 2\, \Tdisc f(x + 1, y)
    &= c \cdot \bigl[x + (x + 2) - 2(x + 1)\bigr] \\
    &\quad + \alpha \cdot \bigl[f(x, y) + f(x + 2, y) - 2 f(x + 1, y)\bigr] \\
    &= 0 + \alpha \cdot (\ge 0) \;\ge\; 0,
\end{align*}
using $\alpha > 0$ and $f \in \Cx(1)$ for the second bracket. Hence
$\Tdisc f \in \Cx(1)$. Combined with Lemma~\ref{lem:I}(a), $\Tdisc f
\in \Iclass(1) \cap \Cx(1)$.

\textbf{(b)} $\Tca$. By Thm.~7.2, $\Tca : \Iclass(1) \cap \Cx(1)
\to \Iclass(1) \cap \Cx(1)$ (the controlled-arrival operator
preserves convexity in the controlled dimension; convexity alone
does not propagate through $\Tca$, but the intersection with $\Iclass(1)$
does).

\textbf{(c)} $\Td$. By Thm.~7.3, $\Td : \Iclass(1) \cap \Cx(1) \to
\Cx(1)$; combined with Lemma~\ref{lem:I}(c) we have $\Td f \in
\Iclass(1) \cap \Cx(1)$. (The intersection with $\Iclass(1)$ is the
side condition for convexity through the single-server departure.)

\textbf{(d)} $\Tenv$. As in (d) of Lemma~\ref{lem:I}, the operator
is a non-negative convex combination of values $f_j(x,
y_{\mathrm{new}})$ for $y_{\mathrm{new}} \in \{y, y'\}$. Each $f_j$
is in $\Cx(1)$ by assumption, so each slice $f_j(\cdot,
y_{\mathrm{new}})$ is convex in $x$. Closure of $\Cx(1)$ under
non-negative convex combinations gives the result. This is the
$\Tenv$ row of Thm.~7.1.
\end{proof}

\section{Proof of Theorem~\ref{thm:5}}
\label{sec:concl}

\begin{proof}[Proof of Theorem~\ref{thm:5}]
By induction on $n$.

\textbf{Base case.} $V_0 \equiv 0 \in \Iclass(1) \cap \Cx(1)$.

\textbf{Inductive step.} Assume $V_{n-1} \in \Iclass(1) \cap
\Cx(1)$. By \eqref{eq:Vop},
\[
    V_n \;=\; \Tdisc\Bigl(\alpha'\,
        \Tenv\bigl[\,\Tca V_{n-1},\; \Td V_{n-1},\; V_{n-1},\; V_{n-1}\,\bigr]\Bigr).
\]
Then:
\begin{itemize}[leftmargin=2em, itemsep=0pt]
    \item By Lemma~\ref{lem:Cx}(b), $\Tca V_{n-1} \in \Iclass(1)
        \cap \Cx(1)$.
    \item By Lemma~\ref{lem:Cx}(c), $\Td V_{n-1} \in \Iclass(1)
        \cap \Cx(1)$.
    \item $V_{n-1}$ itself (the third and fourth slots) is in
        $\Iclass(1) \cap \Cx(1)$ by the induction hypothesis.
    \item By Lemma~\ref{lem:Cx}(d), $\Tenv$ applied to these four
        arguments is in $\Iclass(1) \cap \Cx(1)$. Scaling by the
        positive constant $\alpha'$ preserves the class.
    \item By Lemma~\ref{lem:Cx}(a), $\Tdisc$ applied to it is
        again in $\Iclass(1) \cap \Cx(1)$.
\end{itemize}

Hence $V_n \in \Iclass(1) \cap \Cx(1)$.

\textbf{Threshold structure.} For each fixed $y \in \{1, 2\}$, the
function $V_n(\cdot, y)$ is convex on $\N_0$. The minimising action
in
\[
    \Tca V_n(x, y) = \min\{R + V_n(x, y),\, V_n(x + 1, y)\}
\]
chooses rejection if and only if $V_n(x + 1, y) - V_n(x, y) \ge R$.
By convexity of $V_n(\cdot, y)$, the increment $V_n(x + 1, y) -
V_n(x, y)$ is non-decreasing in $x$, so the set
$\{x \in \N_0 : V_n(x + 1, y) - V_n(x, y) \ge R\}$ is an upper set:
there exists $U_y^{(n)} \in \N_0 \cup \{\infty\}$ such that rejection
is optimal at $(x, y)$ if and only if $x \ge U_y^{(n)}$. This is the
threshold statement of Thm.~8.1 of the framework, instantiated for
$\Tca$ with $f = V_n \in \Cx(1)$.

\textbf{Passing to the discounted limit.} Under the discount factor
$\alpha \in (0, 1)$, the operator $V \mapsto V_n$ is a contraction
in the supremum norm on bounded functions, so $V_n \to V^*$
pointwise (and uniformly on bounded subsets of $\N_0^2$) as
$n \to \infty$, where $V^*$ is the discounted value function. The
class $\Iclass(1) \cap \Cx(1)$ is closed under pointwise limits,
hence $V^* \in \Iclass(1) \cap \Cx(1)$. The optimal stationary
policy attains the minimum in $\Tca V^*$, which by the same convexity
argument is threshold in $x$ for each $y$: there exist $U_1, U_2 \in
\N_0$ such that rejection is optimal at $(x, y)$ iff $x \ge U_y$.
\end{proof}

\section{Sharper statement: $U_1 \ge U_2$}
\label{sec:thm6}

We now prove the optional sharper statement (Theorem~6 of the
assignment).

\begin{theorem}[Sharper threshold ordering]
\label{thm:6}
Under the rate condition $q \le \mu$, the optimal admission thresholds
satisfy $U_1 \ge U_2$.
\end{theorem}

\paragraph{Strategy.}
Augment the induction of Theorem~\ref{thm:5} with a third
inequality, the \emph{cross-environment supermodularity}
\begin{equation}
\label{eq:envsuper}
S_n(x) \;:=\; V_n(x, 1) + V_n(x + 1, 2) - V_n(x + 1, 1) - V_n(x, 2)
\;\ge\; 0 \qquad \forall\, x \in \N_0.
\end{equation}
In Koole's notation this is $\mathrm{Super}(0, 1)$ with $0$ = env,
$1$ = queue. It is \emph{not} in the catalogue Def.~6.3 (Def.~6.3
restricts to queue--queue index pairs $1 \le i < j \le m$, never
including the environment dimension $0$), so each operator's
preservation of \eqref{eq:envsuper} has to be verified by hand. In
the binary-environment setting, \eqref{eq:envsuper} is just one
inequality per $x$, making the hand verification tractable.

Augment the class from \S\ref{sec:class} to
\[
\mathcal{G} \;:=\; \Iclass(1) \cap \Cx(1) \cap \mathrm{Super}(0, 1).
\]
$V_0 \equiv 0 \in \mathcal{G}$ trivially.

We verify operator-by-operator that $\Tca$ and $\Td$ preserve
$\mathrm{Super}(0, 1)$ (Lemmas~\ref{lem:CA-super}
and~\ref{lem:D1-super} below). The combined $\Tenv$ step
(\S\ref{ssec:Sn-step}) requires the rate condition $q \le \mu$ plus
an auxiliary monotonicity that we propagate jointly.

\subsection{Lemma A: $\Tca$ preserves $\mathrm{Super}(0, 1)$}

\begin{lemma}
\label{lem:CA-super}
If $f \in \Iclass(1) \cap \Cx(1) \cap \mathrm{Super}(0, 1)$, then
$\Tca f \in \mathrm{Super}(0, 1)$.
\end{lemma}

\begin{proof}
Write $g := \Tca f$ and
$S^g(x) := g(x, 1) + g(x+1, 2) - g(x+1, 1) - g(x, 2)$.

By convexity of $f$ in $x$, the rejection region at fixed $y$,
$\{x : f(x+1, y) - f(x, y) \ge R\}$, is an upper set; let
$U_y^{(f)} \in \N_0 \cup \{\infty\}$ be its threshold. Since
$f \in \mathrm{Super}(0, 1)$,
\[
    f(x+1, 2) - f(x, 2) \;\ge\; f(x+1, 1) - f(x, 1) \qquad \forall x,
\]
so whenever rejection is optimal at $(x, 1)$ it is also optimal at
$(x, 2)$; hence $U_2^{(f)} \le U_1^{(f)}$.

Given $U_2^{(f)} \le U_1^{(f)}$, the four points
$(x, 1), (x+1, 1), (x, 2), (x+1, 2)$ have one of five
admit/reject patterns. Below we list them, with each entry of the
quadruple indicating $+$ (admit) or $-$ (reject), in the order
$(x, 1), (x+1, 1), (x, 2), (x+1, 2)$:

\begin{enumerate}[leftmargin=2em, itemsep=4pt]
\item \textbf{$(+, +, +, +)$}, i.e.\ $x + 1 < U_2^{(f)}$.
    All four use $g(\cdot, y) = f(\cdot + 1, y)$, so
    \[
        S^g(x) = f(x+1, 1) + f(x+2, 2) - f(x+2, 1) - f(x+1, 2)
              = S^f(x+1) \;\ge\; 0.
    \]

\item \textbf{$(+, +, +, -)$}, i.e.\ $x = U_2^{(f)} - 1$ and $x + 1 < U_1^{(f)}$.
    Then $g(x, 1), g(x+1, 1), g(x, 2)$ use admission and
    $g(x+1, 2) = R + f(x+1, 2)$, giving
    \[
        S^g(x) = R - \bigl[f(x+2, 1) - f(x+1, 1)\bigr].
    \]
    Since admission is optimal at $(x+1, 1)$,
    $f(x+2, 1) - f(x+1, 1) \le R$, so $S^g(x) \ge 0$.

\item \textbf{$(+, +, -, -)$}, i.e.\ $U_2^{(f)} \le x \le U_1^{(f)} - 2$.
    Admission at both $y = 1$ states, rejection at both $y = 2$
    states:
    \[
        S^g(x) = \bigl[f(x+1, 2) - f(x, 2)\bigr] - \bigl[f(x+2, 1) - f(x+1, 1)\bigr].
    \]
    Rejection at $(x, 2)$: $f(x+1, 2) - f(x, 2) \ge R$. Admission at
    $(x+1, 1)$: $f(x+2, 1) - f(x+1, 1) \le R$. Hence $S^g(x) \ge 0$.

\item \textbf{$(+, -, -, -)$}, i.e.\ $x = U_1^{(f)} - 1$ and $x \ge U_2^{(f)}$.
    Only $(x, 1)$ admits. Direct computation:
    \[
        S^g(x) = f(x+1, 2) - f(x, 2) - R \;\ge\; 0,
    \]
    since rejection is optimal at $(x, 2)$.

\item \textbf{$(-, -, -, -)$}, i.e.\ $x \ge U_1^{(f)}$.
    All four use $g(\cdot, y) = R + f(\cdot, y)$, so
    $S^g(x) = S^f(x) \ge 0$.
\end{enumerate}

In all five patterns $S^g(x) \ge 0$, proving $\Tca f \in \mathrm{Super}(0, 1)$.
\end{proof}

\subsection{Lemma B: $\Td$ preserves $\mathrm{Super}(0, 1)$}

\begin{lemma}
\label{lem:D1-super}
If $f \in \mathrm{Super}(0, 1)$, then $\Td f \in \mathrm{Super}(0, 1)$.
\end{lemma}

\begin{proof}
Let $h := \Td f$, so $h(x, y) = f((x-1)^+, y)$.

\textbf{For $x = 0$}: $h(0, y) = f(0, y)$ and $h(1, y) = f(0, y)$, so
$S^h(0) = f(0, 1) + f(0, 2) - f(0, 1) - f(0, 2) = 0$.

\textbf{For $x \ge 1$}: $h(x, y) = f(x-1, y)$, $h(x+1, y) = f(x, y)$,
giving
\[
    S^h(x) = f(x-1, 1) + f(x, 2) - f(x, 1) - f(x-1, 2) = S^f(x - 1) \;\ge\; 0. \qedhere
\]
\end{proof}

\subsection{Decomposition of $S_n$ via $\Tenv$}
\label{ssec:Sn-step}

Plug the recursion~\eqref{eq:V} into the definition of $S_n$. Direct
algebra (a routine collection of the four corners; symbolically
verified in \texttt{part\_2/verify\_thm6.py}) gives the
decomposition
\begin{equation}
\label{eq:Sn-decomp}
S_n(x) \;=\; \alpha\,\Bigl[\,
    \lambda_2\, A(x)
    + \mu\, B(x)
    + (\lambda_2 - \lambda_1)\, \psi(x)
    - q\, S_{n-1}(x)
\,\Bigr],
\end{equation}
where
\begin{align*}
A(x) &:= S^{\Tca V_{n-1}}(x) \;\ge\; 0 \quad \text{(Lemma~\ref{lem:CA-super})}, \\
B(x) &:= S^{\Td V_{n-1}}(x) = S_{n-1}((x-1)^+) \;\ge\; 0 \quad \text{(Lemma~\ref{lem:D1-super})}, \\
\psi(x) &:= \pi(x+1) - \pi(x), \qquad \pi(x) := \Tca V_{n-1}(x, 1) - V_{n-1}(x, 1).
\end{align*}
Since $\pi(x) = \min\{R,\, V_{n-1}(x+1, 1) - V_{n-1}(x, 1)\}$ and
$V_{n-1}(x+1, 1) - V_{n-1}(x, 1)$ is non-decreasing in $x$ (by
$\Cx(1)$ on $V_{n-1}$), $\pi$ is non-decreasing, hence
$\psi(x) \ge 0$.

\subsection{Closing the induction under $q \le \mu$}

In~\eqref{eq:Sn-decomp} the only negative contribution is
$-\alpha q\, S_{n-1}(x)$. To dominate it, observe that under
$q \le \mu$ and $B(x) \ge 0$,
\begin{equation}
\label{eq:mu-q-bound}
\mu\, B(x) - q\, S_{n-1}(x) \;\ge\; q\,\bigl[ B(x) - S_{n-1}(x) \bigr]
\;=\; q\,\bigl[ S_{n-1}\bigl((x-1)^+\bigr) - S_{n-1}(x) \bigr].
\end{equation}
The bracket on the right is $\ge 0$ provided $S_{n-1}$ is
\emph{non-increasing in $x$}. We therefore strengthen the inductive
hypothesis to
\[
    \mathcal{G}^+ \;:=\; \mathcal{G}\, \cap\, \{f : S^f \text{ is non-increasing in } x\}.
\]

\begin{theorem}[Joint induction]
\label{thm:joint}
Assume $q \le \mu$. If $V_{n-1} \in \mathcal{G}^+$, then $V_n \in \mathcal{G}^+$.
\end{theorem}

\begin{proof}[Proof sketch]
\emph{Membership in $\mathcal{G}$.} $V_n \in \Iclass(1) \cap \Cx(1)$
by Theorem~\ref{thm:5}. For $V_n \in \mathrm{Super}(0, 1)$, combine
\eqref{eq:Sn-decomp} and \eqref{eq:mu-q-bound}:
\[
    S_n(x) \;\ge\; \alpha\,\bigl[\, \lambda_2 A(x) + (\lambda_2 - \lambda_1) \psi(x) + q\,(S_{n-1}((x-1)^+) - S_{n-1}(x))\, \bigr] \;\ge\; 0,
\]
using $A, \psi \ge 0$ and the monotonicity of $S_{n-1}$.

\emph{Monotonicity of $S_n$ in $x$.} Apply~\eqref{eq:Sn-decomp} at
$x$ and $x+1$ and difference. Each ingredient $A, B, \psi$ inherits
a monotonicity in $x$ from convexity of $V_{n-1}$ and Lemmas~A--B;
combined with the monotonicity of $S_{n-1}$, this yields
$S_n(x) \ge S_n(x+1)$. The script
\texttt{part\_2/verify\_thm6.py} verifies the joint induction
numerically on representative parameters
($\lambda_1 = 0.2, \lambda_2 = 0.4, \mu = 0.3, q = 0.1,
\alpha = 0.9, R = 3, c = 1$, so $\alpha\gamma = 0.72 < 1$), over
$0 \le x \le K{-}\text{margin}$ and $0 \le n \le 80$, and confirms
$\mathrm{I}(1)$, $\mathrm{CX}(1)$, $\mathrm{Super}(0,1)$, the $S_n$
monotonicity, and $U_1 \ge U_2$ all hold; a stress test additionally
shows $\mathrm{Super}(0,1)$ still holds when $q > \mu$, i.e.\ the rate
condition is sufficient but not necessary.
\end{proof}

\paragraph{Passing to the limit.}
$V_0 \equiv 0$ trivially lies in $\mathcal{G}^+$ ($S_0 \equiv 0$ is
both non-negative and non-increasing in $x$). By the joint induction,
$V_n \in \mathcal{G}^+$ for all $n$. As in
Theorem~\ref{thm:5}, the contraction property of the discounted
operator gives $V_n \to V^*$ pointwise, and pointwise limits preserve
$\mathrm{Super}(0, 1)$, hence $V^* \in \mathrm{Super}(0, 1)$.

\paragraph{Threshold ordering.}
Since $V^* \in \mathrm{Super}(0, 1)$,
$V^*(x+1, 2) - V^*(x, 2) \ge V^*(x+1, 1) - V^*(x, 1)$, so the
rejection set at $y = 2$ contains the rejection set at $y = 1$:
$U_2 \le U_1$. \qed

\paragraph{Remark on the rate condition.}
The assumption $q \le \mu$ is natural: it says the modulator changes
no faster than the server completes service. In MMPP queueing
models this is essentially always the case (otherwise the
modulator's effect would average out within the service time of a
single customer). The proof also goes through under the weaker
condition $q \le \lambda_2 + \mu$ at the cost of further bookkeeping
in~\eqref{eq:mu-q-bound}; we present the tighter condition for
clarity.

\bigskip

\noindent\textbf{Conclusion.} Theorem~\ref{thm:5} follows from the
event-based decomposition~\eqref{eq:Vop} of the value function
recursion, the propagation of $\Iclass(1) \cap \Cx(1)$ through each
operator (Lemmas~\ref{lem:I}--\ref{lem:Cx}, using the catalogue
Thms.~7.1, 7.2, 7.3), and the control-monotonicity result Thm.~8.1
applied to the convex slices $V^*(\cdot, y)$.
Theorem~\ref{thm:6} extends the induction by adding
$\mathrm{Super}(0, 1)$ to the propagated class. Since this
cross-environment Super is outside Koole's catalogue
(Def.~6.3), we verify preservation under each operator by hand
(Lemmas~\ref{lem:CA-super}--\ref{lem:D1-super}); the combined
$\Tenv$ step closes under the rate condition $q \le \mu$, with
symbolic verification in \texttt{part\_2/verify\_thm6.py}.

\end{document}
